\scanpage{11}
\mspage{9}
Ce qui donne un espoir que le complexe de DR à p.d. de $\unsure{\mathscr X|X|\mathcal X}_\ast$ \inserted{à coeff ds $k=\mathbb Z$ disons} (même dans le cas $k=\mathbb Q$ \unsure{\}|)} qu'il contient) donne \unsure{une|bien} information homotopique précise sur $\mathscr X_\ast$, c'est l'observation qu'en fait, pour un \deleted{module gradué} $k$-module $M$ quelconque, on récupère $M$ à isom canonique près par la connaissance d'un quelconque des tronqués $\tau_q(M\otimes_k S^\bullet)$ (quelque grand que soit $q$\ldots) \inserted{\unsure{que j'appelle|qu'on appelle} la \uline{$q$ième ombre de $M$}.}, et de $=$ tt accouplement $M\otimes N\unsure{\text{---}|\to}P$ est connu quand on connaît, pour $q$ assez grand, l'accouplement correspondant $\tau_q(M_S)\otimes\tau_q(N_{\unsure{S|J}})\unsure{\text{---}|\to}\tau_q(P_S)$. Donc une \inserted{$k$-}algèbre graduée $H^\bullet$ à degrés $\geq 0$ \deleted{est} est connue à isom canonique près quand on connaît la $S$-algèbre bigraduée dont les composantes \inserted{de degré} ``extérieur'' \inserted{$q$} sont les $\tau_q(H^q\otimes_k S)$ \ldots. \deleted{Ainsi} Cette ``théorie des ombres'' étant \inserted{supposée} acquise, on voit que la \uline{connaissance} de $C^{\bullet\bullet}_{\mathrm{DRpd}}(\mathscr X_\ast,k)$ \inserted{(en tant que $S$-algèbre diff. bigraduée)} implique celle de l'algèbre graduée $H^{\unsure{\ast|\bullet}}(\mathscr X_\ast,k)$ -- et en raffinant un peu, on \inserted{voit qu'on} trouve \unsure{que|$=$} $R^{\bullet}\unsure{\Gamma|\mathscr T}(\mathscr X_\ast,k)$ comme élément de la cat. dérivée $D(k)$, \uline{plus} l'accouple%
% archive sheet number: 10
